\subsection{Conjugacy Theorems}

THEOREM 2. --- Let B be a semisimple ring and Z its center; suppose that B is a finitely generated Z-module. Let A be a right Artinian ring, and let f and g be ring homomorphisms from B to A; write \(f_{Z}\) and \(g_{Z}\) for the restrictions of f and g to Z. The following properties are equivalent:

\begin{enumerate}
\def\labelenumi{(\roman{enumi})}
\item
  There exists an inner automorphism \(\theta\) of A such that \(g = \theta \circ f\) .
\item
  There exists an inner automorphism \(\theta\) of A such that \(g_{\mathbb{Z}} = \theta \circ f_{\mathbb{Z}}\) .
\end{enumerate}

Since the ring A is right Artinian, \(A_{d}\) is a right A-module of finite length (VIII, p.~6, Theorem 1). Hence \(A^{f}\) and \(A^{g}\) are (B, A)-bimodules of finite length. By Lemma 1 (VIII, p.~254), assertion (i) means that \(A^{f}\) and \(A^{g}\) are isomorphic (B, A)-bimodules and assertion (ii) that they are isomorphic (Z, A)-bimodules. The equivalence of (i) and (ii) therefore follows from Lemma 2 (VIII, p.~254).

COROLLARY. --- Let A and B be algebras over the field K. Suppose that B is central simple and of finite degree and that A is right Artinian. Let f and g be K-algebra homomorphisms from B to A. There exists an inner automorphism \(\theta\) of A such that \(g = \theta \circ f\) .

In the notation of Theorem 2, we have \(\mathrm{Z} = \mathrm{K}\) and therefore \(f_{\mathrm{Z}} = 1_{\mathrm{Z}} = g_{\mathrm{Z}}\) .

THEOREM 3 (Skolem-Noether). --- Let A and B be simple K-algebras and Z(A) and Z(B) their centers. Suppose that the algebra B has finite degree over K and that the algebra \(Z(A) \otimes_{K} Z(B)\) is a field (which is, in particular, the case when A or B is central). Let f and g be K-algebra homomorphisms from B to A. There exists an inner automorphism \(\theta\) of A such that \(g = \theta \circ f\) .

By Lemma 1 of VIII, p.~254, it suffices to prove that the (B,A)-bimodules \(A^{f}\) and \(A^{g}\) are isomorphic. Now, we can view \(A^{f}\) and \(A^{g}\) as left modules over the algebra \(C = B \otimes_{K} A^{o}\) , which is simple by Proposition 7 of VIII, p.~221. As right A-modules, \(A^{f}\) and \(A^{g}\) are isomorphic to \(A_{d}\) , hence have finite length because the ring A is simple (VIII, p.~121, Corollary 1). A fortiori, \(A^{f}\) and \(A^{g}\) are C-modules of finite length. Let S be a simple C-module; there exist strictly positive integers m and n such that \(A^{f}\) is isomorphic to \(S^{m}\) and \(A^{g}\) to \(S^{n}\) . The right A-module S therefore has nonzero finite length. Since the right A-modules underlying \(A^{f}\) and \(A^{g}\) are isomorphic, they have the same length; we therefore have m = n, so that the C-modules \(A^{f}\) and \(A^{g}\) are isomorphic.

COROLLARY 1. --- Let A be a central simple algebra over K, and let L be an extension of K of finite degree. If f and g are K-algebra homomorphisms from L to A, then there exists an inner automorphism \(\theta\) of A such that \(g = \theta \circ f\) .

COROLLARY 2. --- Let A be a central simple algebra over K, and let L be a subalgebra of A that is a field. Every K-algebra homomorphism from L to A extends to an inner automorphism of A.

COROLLARY 3. --- Let D be a field of finite degree over K with center K. Every element of D is algebraic over K. Let x and y be elements of D; there exists an element a of \(D^{*}\) such that \(y = axa^{-1}\) if and only if x and y have the same minimal polynomial over K.

The first assertion follows from Corollary 1 of V, §3, No.~1, p.~17.

Suppose that there exists an element \(a\) of \(\mathrm{D}^*\) such that \(y = axa^{-1}\) ; for every polynomial \(\mathrm{P}\) of \(\mathrm{K}[\mathrm{X}]\) , we have \(\mathrm{P}(y) = a\mathrm{P}(x)a^{-1}\) , and, in particular, we have \(\mathrm{P}(x) = 0\) if and and only if \(\mathrm{P}(y) = 0\) . Consequently, \(x\) and \(y\) have the same minimal polynomial over \(\mathrm{K}\) (V, §3, No.~1, p.~16, Theorem 1).

Conversely, suppose that x and y have the same minimal polynomial. By loc. cit., there exists a K-isomorphism u from K{[}x{]} to K{[}y{]} such that \(u(x) = y\) , and K{[}x{]} is a field. By Corollary 2, u extends to an inner automorphism \(\theta : z \mapsto aza^{-1}\) of D, and we therefore have \(y = \theta(x) = axa^{-1}\) .

PROPOSITION 1. --- Let A be a central simple algebra of finite degree over K. Let B be a K-algebra, and let f and g be algebra homomorphisms from B to A. The following properties are equivalent:

\begin{enumerate}
\def\labelenumi{(\roman{enumi})}
\item
  There exists an inner automorphism \(\theta\) of A such that \(g = \theta \circ f\) .
\item
  As left B-modules, \(\mathrm{A}^f\) and \(\mathrm{A}^g\) are isomorphic.
\end{enumerate}

By Lemma 1 (VIII, p.~254), property (i) is equivalent to the fact that \(A^{f}\) and \(A^{g}\) are isomorphic as (B, A)-bimodules. Since A is finite-dimensional over K, \(A^{f}\) and \(A^{g}\) are B-modules of finite length. Since the center of A is equal to K, the equivalence of (i) and (ii) follows from Lemma 2 of VIII, p.~254 applied to the \((\mathrm{A}^{\circ}, \mathrm{B}^{\circ})\) -bimodules \(A^{f}\) and \(A^{g}\) .

